\documentclass[10pt,a4paper]{amsart}
\usepackage[margin=19mm]{geometry}
\usepackage{amsmath,amssymb,mathtools}
\usepackage[scr=rsfs,cal=euler]{mathalfa}
\usepackage{xfrac}
\usepackage[unicode,hidelinks,bookmarks=false]{hyperref}
\newcommand{\ie}{i.e.\ }
\newcommand{\cf}{cf.\ }
\newcommand{\resp}{respectively\ }
\newcommand{\Subsection}{\subsection}
\providecommand{\nobreakdash}{\nobreak\hbox{-}}
\setlength{\emergencystretch}{2em}
\multlinegap0pt
\usepackage{amssymb}
\usepackage{mathtools}
\newtheorem{lemma}{Lemma}
\newcommand{\R}{\mathbb{R}}
\newcommand{\grad}{\nabla}
\newcommand{\lapl}{\Delta}
\title{Scaling Inequalities and Limits for Clamped Plate Eigenvalues on Geodesic Disks}
\author{Scott Harman}
\date{Source: arXiv:2608.29498v1 (CC BY 4.0)}
\begin{document}
\maketitle

\noindent\emph{Source and licence.} Scaling Inequalities and Limits for Clamped Plate Eigenvalues on Geodesic Disks. Version arXiv:2608.29498v1 (CC BY 4.0), distributed by arXiv under the Creative Commons Attribution 4.0 International licence, http://creativecommons.org/licenses/by/4.0/, as recorded in the arXiv licence metadata for this exact version and shown on the abstract page https://arxiv.org/abs/2608.29498v1. Full licence notice: \texttt{licenses/arxiv-2608.29498-CC-BY-4.0.txt}.\par\smallskip

\noindent\textit{Editorial scope.} Counted unit: the article's lemma lem:ptwisevanfourthorder--numer (source0.tex lines 701--716). The article's complete original proof of it is delivered with it (source0.tex lines 766--805, one \texttt{proof} environment of 40 lines, which the article places away from the statement). No mathematics is written, repaired or re-derived: the statement and the proof are the article's own lines, in the article's own notation, and no retained line is substituted or re-flowed.  Retained uncounted as the source-local context the proof needs: lines 679--699.  Nothing was omitted from any retained span, so the package carries no omission record.  No retained line contains a citation, so the wrapper carries no bibliography and no bibliography entry.  No retained line contains a figure environment, an image-inclusion command or drawing code, so no image is delivered and the package declares no asset dependency.  Editorial formatting only, in addition: an AMS \texttt{amsart} preamble that restates the article's own declarations and macros used by the retained lines, namely the environments \texttt{lemma} on separate counters, together with the standard packages those lines need.  The retained environments print in this document as Lemma 1 in order of appearance, whereas the article prints the same items with the numbering of its own full document.  The complete native source of the article as distributed in the arXiv e-print for this exact version is bundled unchanged as \texttt{sources/208856/source0.tex}.
\begin{lemma}[Pointwise vanishing of the denominator]\label{lem:ptwisevanfourthorder}

Let $\Omega \subset \R^2$ be a bounded domain with smooth boundary. For each $R > 0$, let $w_R:\Omega \to (0,1]$ be a measurable function with $\text{ess inf } w_R > 0$. Fix $\tau \geq 0$ and $\rho \geq 0$. Let $\Gamma_k(w_R,\tau)$ and $\Lambda_k(w_R,\rho)$ be given by the variational characterizations
$$
\Gamma_k(w_R,\tau)
=
\min_{\mathcal L}\ \max_{\substack{0\neq v\in \mathcal L}}
\frac{\displaystyle \int_\Omega (\Delta v)^2\,dA + \tau \int_\Omega |\nabla v|^2\,dA}
{\displaystyle \int_\Omega v^2\,w_R\,dA},
$$
$$
\Lambda_k(w_R,\rho)
=
\min_{\mathcal L}\ \max_{\substack{0\neq v\in \mathcal L}}
\frac{\displaystyle \int_\Omega (\Delta v)^2\,dA + \rho \int_\Omega v^2\,dA}
{\displaystyle \int_\Omega |\nabla v|^2\,w_R\,dA},
$$
where $k \geq 1$ and $\mathcal L$ varies over $k$-dimensional subspaces of $H_0^2(\Omega)$. 

If $\lim_{R\to\infty} w_R = 0$ a.e., then $\lim_{R\to\infty} \Gamma_k(w_R,\tau)=\infty$ $\lim_{R\to\infty} \Lambda_k(w_R,\rho)=\infty$.
\end{lemma}

\begin{lemma}[Pointwise explosion of the numerator]\label{lem:ptwisevanfourthorder--numer}
Let $\Omega \subset \R^2$ be a bounded domain with smooth boundary. For each $R > 0$, let $w_R: \Omega\to (0,1]$ be a measurable function with $\text{ess inf } w_R > 0$. Fix $\tau \geq 0$ and $\rho \geq 0$. Let $\Gamma_k(w_R, \tau)$ and $\Lambda_k(w_R, \rho)$ be given by the variational characterizations
$$
\Gamma_k(w_R, \tau)
= \min_{\mathcal L}\ \max_{\substack{0\neq v\in\mathcal L}}
\frac{\displaystyle \int_{\Omega}\left((\Delta v)^2 + \tau |\nabla v|^2\right)w_R^{-1}\,dA}{\displaystyle \int_{\Omega} v^2\,dA},
$$
$$
\Lambda_k(w_R, \rho)
= \min_{\mathcal L}\ \max_{\substack{0\neq v\in\mathcal L}}
\frac{\displaystyle \int_{\Omega}\left((\Delta v)^2 + \rho v^2)\right)w_R^{-1}\,dA}{\displaystyle \int_{\Omega} |\nabla v|^2\,dA},
$$
where $k \geq 1$ and $\mathcal L$ varies over $k$-dimensional subspaces of $H_0^2(\Omega)$. 

If $\lim_{R\to\infty} w_R = 0$ a.e., then $\lim_{R\to\infty} \Gamma_k(w_R,\tau)=\infty$ and $\lim_{R\to\infty} \Lambda_k(w_R,\rho)=\infty$.
\end{lemma}

\begin{proof}[Proof of Lemma \ref{lem:ptwisevanfourthorder--numer}] This time, the variational characterizations of the eigenvalues associated to the weight $m_R$ are
$$
\Gamma_k(m_R, \tau)
= \min_{\mathcal L}\ \max_{\substack{0\neq v\in\mathcal L}}
\frac{\displaystyle \int_{\Omega}\left((\Delta v)^2 + \tau |\nabla v|^2\right)m_R^{-1}\,dA}{\displaystyle \int_{\Omega} v^2\,dA},
$$
$$
\Lambda_k(m_R, \rho)
= \min_{\mathcal L}\ \max_{\substack{0\neq v\in\mathcal L}}
\frac{\displaystyle \int_{\Omega}\left((\Delta v)^2 + \rho v^2\right)m_R^{-1}\,dA}{\displaystyle \int_{\Omega} |\nabla v|^2\,dA},
$$
where $\mathcal{L}$ varies over $k$-dimensional subspaces of $H_0^2(\Omega)$ and all terms are nonnegative since $\tau, \rho \geq 0$. Clearly, $\Gamma_k(w_R,\tau) \geq \Gamma_k(m_R,\tau)$ and $\Lambda_k(w_R,\rho) \geq \Lambda_k(m_R,\rho)$.  Inspection of the variational characterizations also shows that
\begin{equation}\label{lemmaweak1num}
\delta^{-1}\Gamma_k(1, \tau) \geq \Gamma_k(m_R, \tau) \geq \Gamma_k(1, \tau),
\end{equation}
\begin{equation}\label{lemmaweak2num}
\delta^{-1}\Lambda_k(1, \rho) \geq \Lambda_k(m_R, \rho) \geq \Lambda_k(1, \rho).
\end{equation}
Our goal is to show
\begin{equation}\label{eq:weakgoalexpl1}
\liminf_{R \to \infty} \Gamma_k(m_R,\tau)\geq \delta^{-1}\Gamma_1(1,\tau),
\end{equation}
\begin{equation}\label{eq:weakgoalexpl2}
\liminf_{R \to \infty} \Lambda_k(m_R,\rho)\geq \delta^{-1}\Lambda_1(1,\rho),
\end{equation}
after which the lemma follows by letting $\delta \to 0$.

Following along with the proof of Lemma \ref{lem:ptwisevanfourthorder}, we generate weak eigenfunctions $f_R$ for $\Gamma_k(m_R, \tau)$ and $g_R$ for $\Lambda_k(m_R,\rho)$ that satisfy
$$
\int_\Omega \frac{\lapl \varphi \lapl f_R + \tau \grad \varphi \cdot \grad f_R}{m_R}\, dA
= \Gamma_k(m_R, \tau) \int_\Omega \varphi f_R \, dA,
$$
$$
\int_\Omega \frac{\lapl \varphi \lapl g_R + \rho \varphi g_R}{m_R}\, dA
= \Lambda_k(m_R, \rho) \int_\Omega \grad \varphi \cdot \grad g_R \, dA.
$$
We normalize by $\int_\Omega f_R^2 = 1$ and $\int_\Omega |\grad g_R|^2 = 1$ for each $R$, and in both cases, we generate families of functions uniformly bounded in $H_0^2(\Omega)$ by choosing $\varphi = f_R$ and $\varphi = g_R$ and using \eqref{lemmaweak1num} and \eqref{lemmaweak2num}. The Rellich-Kondrachov theorem yields two functions $f,g \in H_0^2(\Omega)$ so that $f_R \to f$ weakly in $H^2$ and strongly in $L^2$, and $g_R \to g$ weakly in $H^2$ and strongly in $H^1$.

The rest of the proof is identical to that of Lemma \ref{lem:ptwisevanfourthorder}, with the single modification that we use $m_R^{-1}\to \delta^{-1}$ a.e. in place of $m_R \to \delta$ a.e., as $R \to \infty$, hence arriving at \eqref{eq:weakgoalexpl1} and \eqref{eq:weakgoalexpl2}. Letting $\delta \to 0$, we conclude that the eigenvalues of the weighted operators tend to $\infty$ as $R \to \infty$.
\end{proof}

\end{document}
